\documentclass[12pt]{article}
\usepackage[utf8]{inputenc}
\usepackage[T1]{fontenc}
\usepackage{hyperref}
\title{Bimin Language Resources for K-12 Curriculum Development}
\author{Compiled by Claude Sonnet 5.5 (claude-sonnet-5-5) for Andr\'as Kornai}
\date{}
\begin{document}
\maketitle

\section{Introduction}
Bimin (ISO 639-3: bhl; dialects Bim Weng and Nimtew Weng) is a Mountain Ok language of Papua New Guinea, spoken between the Murray and Strickland Rivers in Sandaun and Western Provinces. Search results give about 2,250--2,300 speakers (2003) and 4,300 for the ethnic population, lower than other figures. Little published material was found.

\section{Primary Educational Resources}
No primers, textbooks, or school programmes in Bimin were found in search results.

\section{Linguistic and Reference Materials}
\begin{itemize}
\item Wikipedia, Bimin language: \url{https://en.wikipedia.org/wiki/Bimin_language}
\item Endangered Languages Project entry for Bimin: \url{https://endangeredlanguages.com/lang/10692}
\item Numeral system data for Bimin (a base-27 body-part counting system): \url{https://lingweb.eva.mpg.de/channumerals/Bimin.htm}
\end{itemize}
No dictionary or grammar was found.

\section{Religious and Community Resources}
\begin{itemize}
\item Joshua Project, Bimin people group page, listing a New Testament, audio recordings and the Jesus Film: \url{https://joshuaproject.net/people_groups/10853}
\end{itemize}

\section{Digital and Audio Resources}
\begin{itemize}
\item Omniglot page on Bimin (Bim Weng), which mentions partial Bible translation: \url{https://omniglot.com/writing/bimin.htm}
\end{itemize}

\section{Curriculum Development Context}
Sources disagree on the extent of scripture translation and on speaker numbers, and no educational materials were found. A native-language curriculum would need new documentation work with the community, building on any scripture translation and orthography already in use.

\end{document}